\documentclass[landscape,parskip=half,7pt]{scrartcl}
\usepackage{amsmath,amsfonts,mathtools,amsthm}
\usepackage[landscape, left=1cm, right=1cm, top=1cm, bottom=1cm]{geometry}
\usepackage[backend=biber]{biblatex}
\addbibresource{../Bibo/Experts\string in\string Teams.bib}

\newtheorem{theorem}{Theorem}
\DeclareMathOperator{\argmin}{argmin}

\makeatletter
\newcommand{\stapel}[1]{%					%substack mit align-Möglichkeit (Summe i,j&=1 \\ k&=0 bis n) in Subscript
	\begingroup
	\protect
	\vcenter{%
		\Let@ \restore@math@cr \default@tag
		\baselineskip\fontdimen10 \scriptfont\tw@
		\advance\baselineskip\fontdimen12 \scriptfont\tw@
		\lineskip\thr@@\fontdimen8 \scriptfont\thr@@
		\lineskiplimit\lineskip
		\ialign{\hfil$\m@th\scriptstyle##$&$\m@th\scriptstyle{}##$\hfil\crcr#1\crcr}%
	}%
	\endgroup
}
\makeatother

\begin{document}
	\section{version 1: Linear $pQ$ with $\gamma$ on $T$}
	\begin{align*}
		&\min E(Q,T)=\frac12\int\limits_{0}^{t_F}\int\limits_\Omega p(t)Q(x,t)\,\mathrm{d}x\mathrm dt+\frac\gamma2\int\limits_\Omega |T(x,t_F)-T_\text{fin}(x)|^2\mathrm dx \\
		&\begin{cases}\begin{aligned}
			\rho c\dot T(x,t)&=\nabla\cdot\left(k(x)\nabla T(x,t)\right)+Q(x,t) && \text{in}\ \Omega\times(0,t_F) \\
			-k(x)\partial_\nu T(x,t)&=h(x)\left(T(x,t)-T_\text{out}(x,t)\right) && \text{on}\ \partial\Omega\times[0,t_F]\\
			T(x,0)&=T_\text{ini}(x) &&
		\end{aligned}\end{cases}\tag{SE}\\
		& Q\in U_\text{ad}=\{Q\in L^2(\Omega\times(0,t_F))\,|\,0\leq Q(x,t)\leq Q_\text{max}(x)\ \text{a.e.}\}\\
		%&\text{optional}:\ T(x,t_F)\geq T_\text{fin}(x) (\text{optional if second part of E should be strict})
	\end{align*}
	where
	$T(x,t)$ temperature in K (Henning: I will use °C, but written K is better), $Q(x,t)$ heating output in $\frac W {m^3}$, $p(t)$ price, $\rho$ density in $\frac{kg}{m^3}$, $c$ heat capacity in $\frac{J}{kg\cdot K}$, $k(x)$ heat transfer coefficient in $\frac W{m\cdot K}$, $h(x)$ convective heat transfer coefficient in $\frac W{m^2K}$, $0$ morning start time

	Remember that $\partial_\nu T$ is the outer normal derivative with respect to $x$ and any initial data (or end data) would have to satisfy the PDE.

	\paragraph{SE and WSE}
	In order to solve the (SE) we will use a more basic method, in which we seperate between the spacial and temporal derivatives and first solve the PDE in space using approximations of the time derivatives and then the ODE in time or we try to decouple both variables as much as possible to get a linear system in block structure. This is done by Gridap automatically. Thus, we need the weak formulation only in space (so a semi-weak formulation of the whole system). Here we will give it in its weak form remarking that every time integral of the (WSE) for the (SWF) would be replaced by "for all $t\in(0,t_F)$", because computing the derivative of the energy functional is easier this way without loosing track of the $t$, and get with the test space $H^1_{\Gamma}(\Omega\times(0,t_F))$ where $\Gamma=\Omega\times\{0\}$
	\begin{align*}
		\int\limits_0^{t_F}\int\limits_\Omega \rho c\dot T(x,t)\phi(x,t)\,\mathrm dx\mathrm dt&=\int\limits_0^{t_F}\int\limits_\Omega \nabla\cdot(k(x)\nabla T(x,t)) \phi(x,t)+Q(x,t)\phi(x,t)\,\mathrm dx\mathrm dt\\
		&=\int\limits_0^{t_F}\int\limits_\Omega -k(x)\nabla T(x,t)\cdot\nabla\phi(x,t)+Q(x,t)\phi(x,t)\,\mathrm dx\mathrm dt+\int\limits_0^{t_F}\int\limits_{\partial\Omega} k(x)\partial_\nu T(y,t) \phi(x,t)\,\mathrm dy\mathrm dt \\
		&=\int\limits_0^{t_F}\int\limits_\Omega -k(x)\nabla T(x,t)\cdot\nabla\phi(x,t)+Q(x,t)\phi(x,t)\,\mathrm dx\mathrm dt-\int\limits_0^{t_F}\int\limits_{\partial\Omega} h(y)(T(y,t)-T_\text{out}(y,t)) \phi(x,t)\,\mathrm dy\mathrm dt
	\end{align*}
	giving us the weak formulation
	\begin{align*}
		\int\limits_0^{t_F}\int\limits_\Omega \rho c\dot T(x,t)\phi(x,t)\,\mathrm dx\mathrm dt=\int\limits_0^{t_F}\int\limits_\Omega -k(x)\nabla T(x,t)\cdot\nabla \phi(x,t)+Q(x,t)\phi(x,t)\,\mathrm dx\mathrm dt-\int\limits_0^{t_F}\int\limits_{\partial\Omega} h(y)\left(T(y,t)-T_\text{out}(y)\right)\phi(y,t)\,\mathrm dy\mathrm dt \tag{WSE}
	\end{align*}
	and thus
	\begin{align*}
		&a_\text{SE}(T,\phi)=\int\limits_0^{t_F}\int\limits_\Omega \rho c\dot T(x,t)\phi(x,t)+k(x)\nabla T(x,t)\cdot\nabla \phi(x,t)\,\mathrm dx\mathrm dt+\int\limits_0^{t_F}\int\limits_{\partial\Omega} h(y)T(y,t)\phi(y,t)\,\mathrm dy\mathrm dt,\qquad\qquad\\
		&l_\text{SE}(\phi)=\int\limits_0^{t_F}\int\limits_\Omega Q(x,t)\phi(x,t)\,\mathrm dx\mathrm dt+\int\limits_0^{t_F}\int\limits_{\partial\Omega} h(y)T_\text{out}(y)\phi(y,t)\,\mathrm dy\mathrm dt
	\end{align*}
	By Lax-Milgram there exists a unique solution $T(x,t)$ and thus we can define a solution operator $S:L^2(\Omega\times(0,t_F))\to H^1(\Omega\times(0,t_F))$, $S(Q)=T$.

	Henning: remember that each boundary integral could be rewritten with the Trace operator, might be helpful for the next part. When writing the (WSE) the open interval for $∀ t\in(0,t_F)$ is important.

	\paragraph{Lagrangian + AE}
	From the optimisation functional and the (WSE) we get the Lagrangian as
	\begin{align*}
		L(T,Q,W)=\frac12\int\limits_{0}^{t_F}\int\limits_\Omega p(t)Q(x,t)\,\mathrm{d}x\mathrm dt+\frac\gamma2\int\limits_\Omega |T(x,t_F)-T_\text{fin}(x)|^2\mathrm dx+\int\limits_0^{t_F}\int\limits_\Omega \rho c\dot T(x,t)W(x,t)+k(x)\nabla T(x,t)\cdot\nabla W(x,t)-Q(x,t)W(x,t)\,\mathrm dx\mathrm dt+\int\limits_0^{t_F}\int\limits_{\partial\Omega} h(y)\left(T(y,t)-T_\text{out}(y)\right)W(y,t)\,\mathrm dy\mathrm dt
	\end{align*}
	which gives us the weak formulation of the adjoint (AE) as (notice that the second term contains $\dot\psi$ which in the (WSE) would be zero and thus needs to be solved differently: in that case we are actually only looking for the derivative of the $x$ dependent $T$, so that the derivative $\partial_{T(x)}T(x,t)$ does not give 1 but a time dependent function)
	\begin{align*}
		L_T(T,Q,W)\psi=\gamma\int\limits_\Omega \left(T(x,t_F)-T_\text{fin}(x)\right)\psi(x,t_F)\mathrm dx+\int\limits_0^{t_F}\int\limits_\Omega \rho c\dot \psi(x,t)W(x,t)+k(x)\nabla \psi(x,t)\cdot\nabla W(x,t)\,\mathrm dx\mathrm dt+\int\limits_0^{t_F}\int\limits_{\partial\Omega} h(y)\psi(y,t)W(y,t)\,\mathrm dy\mathrm dt
	\end{align*}
	By Green for the $t$ variable (integration by parts) we get
	\begin{align*}
		\int\limits_0^{t_F}\int\limits_\Omega \rho c\dot \psi(x,t)W(x,t)\,\mathrm dx\mathrm dt=-\int\limits_0^{t_F}\int\limits_\Omega \rho c\psi(x,t)\dot W(x,t)\,\mathrm dx\mathrm dt+\left[\int\limits_\Omega \rho c\psi(x,t)W(x,t)\,\mathrm dx\right]_0^{t_F}=-\int\limits_0^{t_F}\int\limits_\Omega \rho c\psi(x,t)\dot W(x,t)\,\mathrm dx\mathrm dt+\int\limits_\Omega \rho c\psi(x,t_F)W(x,t_F)\,\mathrm dx
	\end{align*}
	when assuming $W(x,0)=0$ a.e. (because $W$ has to be able to operate as test function for (SE)) and thus $\psi\in H^1_\Gamma(\Omega\times(0,t_F))$, which gives the weak version as
	\begin{align*}
		0=\int\limits_\Omega \gamma\left(T(x,t_F)-T_\text{fin}(x)\right)\psi(x,t_F)+\rho c\psi(x,t_F)W(x,t_F)\mathrm dx+\int\limits_0^{t_F}\int\limits_\Omega -\rho c\dot W(x,t)\psi(x,t)+k(x)\nabla \psi(x,t)\cdot\nabla W(x,t)\,\mathrm dx\mathrm dt+\int\limits_0^{t_F}\int\limits_{\partial\Omega} h(y)\psi(y,t)W(y,t)\,\mathrm dy\mathrm dt
	\end{align*}
	or when setting the end condition based on the first integral simply as
	\begin{align*}
		\int\limits_0^{t_F}\int\limits_\Omega \rho c\psi(x,t)\dot W(x,t)-k(x)\nabla \psi(x,t)\cdot\nabla W(x,t)\,\mathrm dx\mathrm dt=-\int\limits_0^{t_F}\int\limits_{\partial\Omega} h(y)\psi(y,t)W(y,t)\,\mathrm dy\mathrm dt \tag{WAE}
	\end{align*}
	Also by Green for any $t\in(0,t_F)$ it is
	\begin{align*}
		\int\limits_\Omega k(x)\nabla\psi(x,t)\cdot\nabla W(x,t)\,\mathrm dx=-\int\limits_\Omega \nabla\cdot\left(k(x)\nabla W(x,t)\right)\psi(x,t)\,\mathrm dx+\int\limits_{\partial\Omega} k(y)\partial_\nu W(y,t)\psi(y,t)\,\mathrm dy
	\end{align*}
	With these it is
	\begin{align*}
		L_T(T,Q,W)\psi=\gamma\int\limits_\Omega \left(T(x,t_F)-T_\text{fin}(x)\right)\psi(x,t)\mathrm dx-\int\limits_0^{t_F}\int\limits_\Omega \rho c\psi(x,t)\dot W(x,t)+\nabla\cdot\left(k(x)\nabla W(x,t)\right)\psi(x,t)\,\mathrm dx\mathrm dt-\int\limits_0^{t_F}\int\limits_{\partial\Omega} h(y)\psi(y,t)W(y,t)-k(y)\partial_\nu W(y,t)\psi(y,t)\,\mathrm dy\mathrm dt+\int\limits_\Omega \rho c\psi(x,t_F)W(x,t_F)\,\mathrm dx
	\end{align*}
	giving the strong version
	\begin{align*}
		\begin{cases}\begin{aligned}
			\rho c\dot W(x,t)&=-\nabla\cdot\left(k(x)\nabla W(x,t)\right)&&\text{in}\ \Omega\times(0,t_F) \\
			k(y)\partial_\nu W(y,t)&=h(y)W(y,t)&&\text{on}\ \partial\Omega\times[0,t_F] \\
			W(x,t_F)&=-\frac{\gamma}{\rho c}\left(T(x,t_F)-T_\text{fin}(x)\right)&&\text{in}\ \Omega
		\end{aligned}\end{cases}\tag{AE}
	\end{align*}
	Thus, we can see that the (AE) has to be solved backwards in time.
	For the (AE) we have
	\begin{align*}
		a_\text{AE}(W,\psi)=\int\limits_0^{t_F}\int\limits_\Omega \rho c\psi(x,t)\dot W(x,t)-k(x)\nabla \psi(x,t)\cdot\nabla W(x,t)\,\mathrm dx\mathrm dt+\int\limits_0^{t_F}\int\limits_{\partial\Omega} h(y)\psi(y,t)W(y,t)\,\mathrm dy\mathrm dt\qquad\qquad l_\text{AE}(\psi)=0
	\end{align*}
	By Lax-Milgram there exists a unique solution $W(x,t)$ and thus a solution operator $S^*:L^2(\Omega)\to H^1(\Omega\times(0,t_F))$, $S^*(T(x,t_F)-T_\text{fin}(x))=W(x,t)$.

	Henning: Actually, we get the (AE) by the variation principle and not by differentiation, which does not fully make sense here (see \ref{sec:Variation}).

	\paragraph{adjoint forward}
	Since the (AE) has to be solved backwards in time, which is not possible with the software we are using (although implementation would be just as is) we introduce $\widetilde W(x,t)=W(x,t_F-t)$. By the (AE) this function satisfies
	\begin{align}
		\begin{cases}\begin{aligned}
				\rho c\dot {\widetilde W}(x,t)&=\nabla\cdot\left(k(x)\nabla \widetilde W(x,t)\right)&&\text{in}\ \Omega\times(0,t_F) \\
				k(y)\partial_\nu \widetilde W(y,t)&=h(y)\widetilde W(y,t)&&\text{on}\ \partial\Omega\times[0,t_F] \\
				\widetilde W(x,0)&=-\frac{\gamma}{\rho c}\left(T(x,t_F)-T_\text{fin}(x)\right)&&\text{in}\ \Omega
		\end{aligned}\end{cases}\tag{AE'}
	\end{align}
	which can now be solved forward in time. The weak version is
	\begin{align}
		\int\limits_0^{t_F}\int\limits_\Omega \rho c\psi(x,t)\dot{\widetilde{W}}(x,t)+k(x)\nabla \psi(x,t)\cdot\nabla \widetilde W(x,t)\,\mathrm dx\mathrm dt=\int\limits_0^{t_F}\int\limits_{\partial\Omega} h(y)\psi(y,t)\widetilde W(y,t)\,\mathrm dy\mathrm dt \tag{WAE}
	\end{align}
	For the (AE') we have
	\begin{align}
		a_\text{AE}(\widetilde W,\psi)=\int\limits_0^{t_F}\int\limits_\Omega \rho c\psi(x,t)\dot {\widetilde W}(x,t)+k(x)\nabla \psi(x,t)\cdot\nabla \widetilde W(x,t)\,\mathrm dx\mathrm dt-\int\limits_0^{t_F}\int\limits_{\partial\Omega} h(y)\psi(y,t)\widetilde W(y,t)\,\mathrm dy\mathrm dt\qquad\qquad l_\text{AE}(\psi)=0
	\end{align}
	which then gives us the original adjoint as $W(x,t)=\widetilde W(x,t_F-t)$.

	\paragraph{Variation} \label{sec:Variation}
	In order to find the adjoint, we introduce a small variation $\delta T$ of $T$ in the respective space and look at the difference this variation makes on $L$ (this is like computing $L(T_1,Q,W)-L(T_2,Q,W)$, remember that the adjoint is a measure for the sensitivity of the functional with respect to $T$).
	\begin{align*}
		\delta_TL(T,Q,W)&=\frac\gamma2\int\limits_\Omega \delta_T|T(x,t_F)-T_\text{fin}(x)|^2\mathrm dx+\int\limits_0^{t_F}\int\limits_\Omega \rho c\dot {\delta T}(x,t)W(x,t)-k(x)\nabla \delta T(x,t)\cdot\nabla W(x,t)\,\mathrm dx\mathrm dt-\int\limits_0^{t_F}\int\limits_{\partial\Omega} h(y)\left(\delta T(y,t)-T_\text{out}(y)\right)W(y,t)\,\mathrm dy\mathrm dt \\
		&=\gamma\int\limits_\Omega (T(x,t_F)-T_\text{fin}(x))\delta T(x,t_F)\mathrm dx-\int\limits_0^{t_F}\int\limits_\Omega \rho c\delta T(x,t)\dot W(x,t)\,\mathrm dx\mathrm dt+\int\limits_\Omega \rho c\delta T(x,t_F)W(x,t_F)\,\mathrm dx
		-\int\limits_0^{t_F}\int\limits_\Omega k(x)\nabla \delta T(x,t)\cdot\nabla W(x,t)\,\mathrm dx\mathrm dt-\int\limits_0^{t_F}\int\limits_{\partial\Omega} h(y)\left(\delta T(y,t)-T_\text{out}(y)\right)W(y,t)\,\mathrm dy\mathrm dt
	\end{align*}
	where we used the same as above and $\delta T(x,0)=0$ because of the initial value condition, which fixes the $T$s there and because the term is linear. Also for the first term we used Taylor giving us $x^2=x_0^2+2x(x-x_0)+o((x-x_0)^2)$, giving us $\delta_x(x-a)^2=(x-a)^2-(x_0-a)^2=2(x-a)(x-x_0)+o(x-x_0)=2(x-a)\delta_x$.

	\paragraph{Solution operator + gradient}
	Now we introduce the solution operator $S$, which for a given $Q$ gives the solution $T$ of the (SE), so $T=S(Q)$. With this we can introduce the reduced energy functional
	\begin{align}
		e(Q)=\int\limits_{0}^{t_F}\int\limits_\Omega p(t)Q(x,t)\,\mathrm{d}x\mathrm dt+\frac\gamma2\int\limits_\Omega |S(Q)(x,t_F)-T_\text{fin}(x)|^2\mathrm dx
	\end{align}
	In order to find the gradient of $e$ we differentiate it and find by noting that $S$ is linear and therefore has derivative $S'(Q)h=Sh$
	\begin{align}
		e'(Q)h&=\int\limits_{0}^{t_F}\int\limits_\Omega p(t)h(x,t)\,\mathrm{d}x\mathrm dt+\gamma\int\limits_\Omega (S(Q)(x,t_F)-T_\text{fin}(x))Sh\mathrm dx\\
		&=\langle p|h\rangle+\langle S(Q)-T_\text{fin} | Sh\rangle|_{t=t_F}=\langle p|h\rangle+\gamma\langle S^*(S(Q)-T_\text{fin}) | h\rangle|_{t=t_F} \equiv p-\rho cW
	\end{align}
	where $S^*$ is the solution operator for (AE) and $W$ is the solution ........

	\section{Linear $pQ$ with $\lambda$ on $Q$}
	\begin{align*}
		&\min E(Q,T)=\frac\gamma2\int\limits_{0}^{t_F}\int\limits_\Omega p(t)Q(x,t)\,\mathrm{d}x\mathrm dt+\frac12\int\limits_\Omega |T(x,t_F)-T_\text{fin}(x)|^2\mathrm dx \\
		&\begin{cases}\begin{aligned}
				\rho c\dot T(x,t)&=\nabla\cdot\left(k(x)\nabla T(x,t)\right)+Q(x,t) && \text{in}\ \Omega\times(0,t_F) \\
				-k(x)\partial_\nu T(x,t)&=h(x)\left(T(x,t)-T_\text{out}(x,t)\right) && \text{on}\ \partial\Omega\times[0,t_F]\\
				T(x,0)&=T_\text{ini}(x) &&
		\end{aligned}\end{cases}\tag{SE}\\
		& Q\in U_\text{ad}=\{Q\in L^2(\Omega\times(0,t_F))\,|\,0\leq Q(x,t)\leq Q_\text{max}(x)\ \text{a.e.}\}
	\end{align*}

	\paragraph{SE and WSE}
	\begin{align*}
		\int\limits_0^{t_F}\int\limits_\Omega \rho c\dot T(x,t)\phi(x,t)\,\mathrm dx\mathrm dt=\int\limits_0^{t_F}\int\limits_\Omega -k(x)\nabla T(x,t)\cdot\nabla \phi(x,t)+Q(x,t)\phi(x,t)\,\mathrm dx\mathrm dt-\int\limits_0^{t_F}\int\limits_{\partial\Omega} h(y)\left(T(y,t)-T_\text{out}(y)\right)\phi(y,t)\,\mathrm dy\mathrm dt \tag{WSE}
	\end{align*}
	and thus
	\begin{align*}
		&a_\text{SE}(T,\phi)=\int\limits_0^{t_F}\int\limits_\Omega \rho c\dot T(x,t)\phi(x,t)+k(x)\nabla T(x,t)\cdot\nabla \phi(x,t)\,\mathrm dx\mathrm dt+\int\limits_0^{t_F}\int\limits_{\partial\Omega} h(y)T(y,t)\phi(y,t)\,\mathrm dy\mathrm dt,\qquad\qquad\\
		&l_\text{SE}(\phi)=\int\limits_0^{t_F}\int\limits_\Omega Q(x,t)\phi(x,t)\,\mathrm dx\mathrm dt+\int\limits_0^{t_F}\int\limits_{\partial\Omega} h(y)T_\text{out}(y)\phi(y,t)\,\mathrm dy\mathrm dt
	\end{align*}
	By Lax-Milgram there exists a unique solution $T(x,t)$ and thus we can define a solution operator $S:L^2(\Omega\times(0,t_F))\to H^1(\Omega\times(0,t_F))$, $S(Q)=T$.

	\paragraph{Lagrangian + AE}
	From the optimisation functional and the (WSE) we get the Lagrangian as
	\begin{align*}
		L(T,Q,W)=\frac\gamma2\int\limits_{0}^{t_F}\int\limits_\Omega p(t)Q(x,t)\,\mathrm{d}x\mathrm dt+\frac12\int\limits_\Omega |T(x,t_F)-T_\text{fin}(x)|^2\mathrm dx+\int\limits_0^{t_F}\int\limits_\Omega \rho c\dot T(x,t)W(x,t)+k(x)\nabla T(x,t)\cdot\nabla W(x,t)-Q(x,t)W(x,t)\,\mathrm dx\mathrm dt+\int\limits_0^{t_F}\int\limits_{\partial\Omega} h(y)\left(T(y,t)-T_\text{out}(y)\right)W(y,t)\,\mathrm dy\mathrm dt
	\end{align*}
	\begin{align*}
		L_T(T,Q,W)\psi=\int\limits_\Omega \left(T(x,t_F)-T_\text{fin}(x)\right)\psi(x,t_F)\mathrm dx+\int\limits_0^{t_F}\int\limits_\Omega \underbrace{\rho c\dot \psi(x,t)W(x,t)}_{\begin{aligned}
				&=-\int\limits_0^{t_F}\int\limits_\Omega \rho c\psi(x,t)\dot W(x,t)\,\mathrm dx\mathrm dt+\left[\int\limits_\Omega \rho c\psi(x,t)W(x,t)\,\mathrm dx\right]_0^{t_F} \\
				&=-\int\limits_0^{t_F}\int\limits_\Omega \rho c\psi(x,t)\dot W(x,t)\,\mathrm dx\mathrm dt+\int\limits_\Omega \rho c\psi(x,t_F)W(x,t_F)\,\mathrm dx
		\end{aligned}}+k(x)\nabla \psi(x,t)\cdot\nabla W(x,t)\,\mathrm dx\mathrm dt+\int\limits_0^{t_F}\int\limits_{\partial\Omega} h(y)\psi(y,t)W(y,t)\,\mathrm dy\mathrm dt
	\end{align*}
	which gives the weak version as
	\begin{align*}
		0=\int\limits_\Omega \left(T(x,t_F)-T_\text{fin}(x)\right)\psi(x,t_F)+\rho c\psi(x,t_F)W(x,t_F)\mathrm dx+\int\limits_0^{t_F}\int\limits_\Omega -\rho c\dot W(x,t)\psi(x,t)+k(x)\nabla \psi(x,t)\cdot\nabla W(x,t)\,\mathrm dx\mathrm dt+\int\limits_0^{t_F}\int\limits_{\partial\Omega} h(y)\psi(y,t)W(y,t)\,\mathrm dy\mathrm dt
	\end{align*}
	or when setting the end condition based on the first integral simply as
	\begin{align*}
		\int\limits_0^{t_F}\int\limits_\Omega \rho c\psi(x,t)\dot W(x,t)-k(x)\nabla \psi(x,t)\cdot\nabla W(x,t)\,\mathrm dx\mathrm dt-\int\limits_0^{t_F}\int\limits_{\partial\Omega} h(y)\psi(y,t)W(y,t)\,\mathrm dy\mathrm dt=0 \tag{WAE}
	\end{align*}
	Also by Green for any $t\in(0,t_F)$ it is
	\begin{align*}
		\int\limits_\Omega k(x)\nabla\psi(x,t)\cdot\nabla W(x,t)\,\mathrm dx=-\int\limits_\Omega \nabla\cdot\left(k(x)\nabla W(x,t)\right)\psi(x,t)\,\mathrm dx+\int\limits_{\partial\Omega} k(y)\partial_\nu W(y,t)\psi(y,t)\,\mathrm dy
	\end{align*}
	With these it is
	\begin{align*}
		L_T(T,Q,W)\psi=\int\limits_\Omega \left(T(x,t_F)-T_\text{fin}(x)\right)\psi(x,t)\mathrm dx-\int\limits_0^{t_F}\int\limits_\Omega \rho c\psi(x,t)\dot W(x,t)+\nabla\cdot\left(k(x)\nabla W(x,t)\right)\psi(x,t)\,\mathrm dx\mathrm dt-\int\limits_0^{t_F}\int\limits_{\partial\Omega} h(y)\psi(y,t)W(y,t)-k(y)\partial_\nu W(y,t)\psi(y,t)\,\mathrm dy\mathrm dt+\int\limits_\Omega \rho c\psi(x,t_F)W(x,t_F)\,\mathrm dx
	\end{align*}
	giving the strong version
	\begin{align*}
		\begin{cases}\begin{aligned}
				\rho c\dot W(x,t)&=-\nabla\cdot\left(k(x)\nabla W(x,t)\right)&&\text{in}\ \Omega\times(0,t_F) \\
				k(y)\partial_\nu W(y,t)&=h(y)W(y,t)&&\text{on}\ \partial\Omega\times[0,t_F] \\
				W(x,t_F)&=-\frac{1}{\rho c}\left(T(x,t_F)-T_\text{fin}(x)\right)&&\text{in}\ \Omega
		\end{aligned}\end{cases}\tag{AE}
	\end{align*}
	Thus, we can see that the (AE) has to be solved backwards in time.
	For the (AE) we have
	\begin{align*}
		a_\text{AE}(W,\psi)=\int\limits_0^{t_F}\int\limits_\Omega \rho c\psi(x,t)\dot W(x,t)-k(x)\nabla \psi(x,t)\cdot\nabla W(x,t)\,\mathrm dx\mathrm dt-\int\limits_0^{t_F}\int\limits_{\partial\Omega} h(y)\psi(y,t)W(y,t)\,\mathrm dy\mathrm dt\qquad\qquad l_\text{AE}(\psi)=0
	\end{align*}
	By Lax-Milgram there exists a unique solution $W(x,t)$ and thus a solution operator $S^*:L^2(\Omega)\to H^1(\Omega\times(0,t_F))$, $S^*(T(x,t_F)-T_\text{fin}(x))=W(x,t)$.

	Henning: Actually, we get the (AE) by the variation principle and not by differentiation, which does not fully make sense here (see \ref{sec:Variation}).

	\paragraph{adjoint forward}
	Since the (AE) has to be solved backwards in time, which is not possible with the software we are using (although implementation would be just as is) we introduce $\widetilde W(x,t)=W(x,t_F-t)$. By the (AE) this function satisfies
	\begin{align}
		\begin{cases}\begin{aligned}
				\rho c\dot {\widetilde W}(x,t)&=\nabla\cdot\left(k(x)\nabla \widetilde W(x,t)\right)&&\text{in}\ \Omega\times(0,t_F) \\
				k(y)\partial_\nu \widetilde W(y,t)&=h(y)\widetilde W(y,t)&&\text{on}\ \partial\Omega\times[0,t_F] \\
				\widetilde W(x,0)&=-\frac{1}{\rho c}\left(T(x,t_F)-T_\text{fin}(x)\right)&&\text{in}\ \Omega
		\end{aligned}\end{cases}\tag{AE'}
	\end{align}
	which can now be solved forward in time. The weak version is
	\begin{align}
		\int\limits_0^{t_F}\int\limits_\Omega \rho c\psi(x,t)\dot{\widetilde{W}}(x,t)+k(x)\nabla \psi(x,t)\cdot\nabla \widetilde W(x,t)\,\mathrm dx\mathrm dt=\int\limits_0^{t_F}\int\limits_{\partial\Omega} h(y)\psi(y,t)\widetilde W(y,t)\,\mathrm dy\mathrm dt \tag{WAE}
	\end{align}
	For the (AE') we have
	\begin{align}
		a_\text{AE}(\widetilde W,\psi)=\int\limits_0^{t_F}\int\limits_\Omega \rho c\psi(x,t)\dot {\widetilde W}(x,t)+k(x)\nabla \psi(x,t)\cdot\nabla \widetilde W(x,t)\,\mathrm dx\mathrm dt-\int\limits_0^{t_F}\int\limits_{\partial\Omega} h(y)\psi(y,t)\widetilde W(y,t)\,\mathrm dy\mathrm dt\qquad\qquad l_\text{AE}(\psi)=0
	\end{align}
	which then gives us the original adjoint as $W(x,t)=\widetilde W(x,t_F-t)$.

	\paragraph{Variation}
	In order to find the adjoint, we introduce a small variation $\delta T$ of $T$ in the respective space and look at the difference this variation makes on $L$ (this is like computing $L(T_1,Q,W)-L(T_2,Q,W)$, remember that the adjoint is a measure for the sensitivity of the functional with respect to $T$).
	\begin{align*}
		\delta_TL(T,Q,W)&=\frac12\int\limits_\Omega \delta_T|T(x,t_F)-T_\text{fin}(x)|^2\mathrm dx+\int\limits_0^{t_F}\int\limits_\Omega \rho c\dot {\delta T}(x,t)W(x,t)-k(x)\nabla \delta T(x,t)\cdot\nabla W(x,t)\,\mathrm dx\mathrm dt-\int\limits_0^{t_F}\int\limits_{\partial\Omega} h(y)\left(\delta T(y,t)-T_\text{out}(y)\right)W(y,t)\,\mathrm dy\mathrm dt \\
		&=\int\limits_\Omega (T(x,t_F)-T_\text{fin}(x))\delta T(x,t_F)\mathrm dx-\int\limits_0^{t_F}\int\limits_\Omega \rho c\delta T(x,t)\dot W(x,t)\,\mathrm dx\mathrm dt+\int\limits_\Omega \rho c\delta T(x,t_F)W(x,t_F)\,\mathrm dx
		-\int\limits_0^{t_F}\int\limits_\Omega k(x)\nabla \delta T(x,t)\cdot\nabla W(x,t)\,\mathrm dx\mathrm dt-\int\limits_0^{t_F}\int\limits_{\partial\Omega} h(y)\left(\delta T(y,t)-T_\text{out}(y)\right)W(y,t)\,\mathrm dy\mathrm dt
	\end{align*}
	where we used the same as above and $\delta T(x,0)=0$ because of the initial value condition, which fixes the $T$s there and because the term is linear. Also for the first term we used Taylor giving us $x^2=x_0^2+2x(x-x_0)+o((x-x_0)^2)$, giving us $\delta_x(x-a)^2=(x-a)^2-(x_0-a)^2=2(x-a)(x-x_0)+o(x-x_0)=2(x-a)\delta_x$.

	\paragraph{Solution operator + gradient}
	Now we introduce the solution operator $S$, which for a given $Q$ gives the solution $T$ of the (SE), so $T=S(Q)$. With this we can introduce the reduced energy functional
	\begin{align}
		e(Q)=\int\limits_{0}^{t_F}\int\limits_\Omega p(t)Q(x,t)\,\mathrm{d}x\mathrm dt+\frac\gamma2\int\limits_\Omega |S(Q)(x,t_F)-T_\text{fin}(x)|^2\mathrm dx
	\end{align}
	In order to find the gradient of $e$ we differentiate it and find by noting that $S$ is linear and therefore has derivative $S'(Q)h=Sh$
	\begin{align}
		e'(Q)h&=\int\limits_{0}^{t_F}\int\limits_\Omega p(t)h(x,t)\,\mathrm{d}x\mathrm dt+\gamma\int\limits_\Omega (S(Q)(x,t_F)-T_\text{fin}(x))Sh\mathrm dx\\
		&=\langle p|h\rangle+\langle S(Q)-T_\text{fin} | Sh\rangle|_{t=t_F}=\langle p|h\rangle+\gamma\langle S^*(S(Q)-T_\text{fin}) | h\rangle|_{t=t_F} \equiv p-\rho cW
	\end{align}
	where $S^*$ is the solution operator for (AE) and $W$ is the solution ........

	\section{version 2: quadratic $pQ$ with $\lambda$ on $T$}
	\begin{align*}
		&\min E(Q,T)=\frac12\int\limits_{0}^{t_F}\int\limits_\Omega (p(t)Q(x,t))^2\,\mathrm{d}x\mathrm dt+\frac\gamma2\int\limits_\Omega |T(x,t_F)-T_\text{fin}(x)|^2\mathrm dx \\
		&\begin{cases}\begin{aligned}
			\rho c\dot T(x,t)&=\nabla\cdot\left(k(x)\nabla T(x,t)\right)+Q(x,t) && \text{in}\ \Omega\times(0,t_F) \\
			-k(x)\partial_\nu T(x,t)&=h(x)\left(T(x,t)-T_\text{out}(x,t)\right) && \text{on}\ \partial\Omega\times[0,t_F]\\
			T(x,0)&=T_\text{ini}(x)&&
		\end{aligned}\end{cases}\tag{SE}\\
		& Q\in U_\text{ad}=\{Q\in L^2(\Omega\times(0,t_F))\,|\,0\leq Q(x,t)\leq Q_\text{max}(x)\ \text{a.e.}\}
	\end{align*}
	In this case mostly everything stays the same.

	\paragraph{SE and WSE}
	We have the weak formulation
	\begin{align*}
		\int\limits_0^{t_F}\int\limits_\Omega \rho c\dot T(x,t)\phi(x,t)\,\mathrm dx\mathrm dt=\int\limits_0^{t_F}\int\limits_\Omega -k(x)\nabla T(x,t)\cdot\nabla \phi(x,t)+Q(x,t)\phi(x,t)\,\mathrm dx\mathrm dt-\int\limits_0^{t_F}\int\limits_{\partial\Omega} h(y)\left(T(y,t)-T_\text{out}(y)\right)\phi(y,t)\,\mathrm dy\mathrm dt \tag{WSE}
	\end{align*}
	and thus
	\begin{align*}
		&a_\text{SE}(T,\phi)=\int\limits_0^{t_F}\int\limits_\Omega \rho c\dot T(x,t)\phi(x,t)+k(x)\nabla T(x,t)\cdot\nabla \phi(x,t)\,\mathrm dx\mathrm dt+\int\limits_0^{t_F}\int\limits_{\partial\Omega} h(y)T(y,t)\phi(y,t)\,\mathrm dy\mathrm dt,\qquad\qquad\\
		&l_\text{SE}(\phi)=\int\limits_0^{t_F}\int\limits_\Omega Q(x,t)\phi(x,t)\,\mathrm dx\mathrm dt+\int\limits_0^{t_F}\int\limits_{\partial\Omega} h(y)T_\text{out}(y)\phi(y,t)\,\mathrm dy\mathrm dt
	\end{align*}
	By Lax-Milgram there exists a unique solution $T(x,t)$ and thus we can define a solution operator $S:L^2(\Omega\times(0,t_F))\to H^1(\Omega\times(0,t_F))$, $S(Q)=T$.

	\paragraph{Lagrangian + AE}
	One difference is the Lagrangian which is
	\begin{align}
		L(T,Q,W)=\frac12\int\limits_{0}^{t_F}\int\limits_\Omega (p(t)Q(x,t))^2\,\mathrm{d}x\mathrm dt+\frac\gamma2\int\limits_\Omega |T(x,t_F)-T_\text{fin}(x)|^2\mathrm dx+\int\limits_0^{t_F}\int\limits_\Omega \rho c\dot T(x,t)W(x,t)-k(x)\nabla T(x,t)\cdot\nabla W(x,t)-Q(x,t)W(x,t)\,\mathrm dx\mathrm dt-\int\limits_0^{t_F}\int\limits_{\partial\Omega} h(y)\left(T(y,t)-T_\text{out}(y)\right)W(y,t)\,\mathrm dy\mathrm dt
	\end{align}


	\paragraph{Solution operator + gradient}
	Now we introduce the solution operator $S$, which for a given $Q$ gives the solution $T$ of the (SE), so $T=S(Q)$. With this we can introduce the reduced energy functional
	\begin{align}
		e(Q)=\frac12\int\limits_{0}^{t_F}\int\limits_\Omega (p(t)Q(x,t))^2\,\mathrm{d}x\mathrm dt+\frac\gamma2\int\limits_\Omega |S(Q)(x,t_F)-T_\text{fin}(x)|^2\mathrm dx
	\end{align}
	In order to find the gradient of $e$ we differentiate it and find by noting that $S$ is linear and therefore has derivative $S'(Q)h=Sh$
	\begin{align*}
		e'(Q)h&=\int\limits_{0}^{t_F}\int\limits_\Omega p^2(t)Q(x,t)h(x,t)\,\mathrm{d}x\mathrm dt+\gamma\int\limits_\Omega (S(Q)(x,t_F)-T_\text{fin}(x))Sh\mathrm dx\\
		&=\langle p^2Q|h\rangle+\gamma\langle S(Q)-T_\text{fin} | Sh\rangle|_{t=t_F}=\langle p^2Q|h\rangle+\langle \gamma S^*(S(Q)-T_\text{fin}) | h\rangle|_{t=t_F} \hat= p^2Q-W
	\end{align*}

	Thus, without box constraints from the condition $e'(Q)=0$ we would get $Q=\frac{1}{p^2}W$ and can use the projection of that onto the box for our case (see bang bang).

	We can also get the gradient by the (VI) which we get by $L_Q=0$ as
	\[0=L_Q(T,Q,W)h=\int\limits_{0}^{t_F}\int\limits_\Omega p^2(t)Q(x,t)h(x,t)\,\mathrm{d}x\mathrm dt-\int\limits_0^{t_F}\int\limits_\Omega W(x,t)h(x,t)\,\mathrm dx\mathrm dt=\langle p^2Q-W,h\rangle_{L^2(\Omega\times(0,t_F))}\]
	This property has to hold for every $h\in Q_\text{ad}-Q$ for the $Q$ from the integral, which gives us the bang-bang result.

	\paragraph{Optimal step size}
	We found out that the gradient of the reduced cost is $\nabla e(Q)=p^2Q-W$, so in the gradient descent algorithm we will go into the negative of that direction and we can actually compute the optimal step size. Consider a given iteration $k$ and the computed control $Q_k$, state $T_k$ and adjoint $W_k$. Also let $v_k=-\nabla e(Q)=-(p^2Q-W)$, $I=[t_0,t_F]$ and define as the step-length-dependent function we want to minimise in the iteration
	$$g(s)=e(Q_k+sv_k)$$
	Then we can rewrite
	\begin{align*}
		g(s)&=e(Q_k+sv_k)\\
		&=\frac12\lVert p(Q_k+sv_k)\rVert^2_{L^2(\Omega\times I)}+\frac\gamma2\lVert S(Q_k+sv_k)-T_{\text{fin}}\rVert^2_{L^2(\Omega)}\big|_{t=t_F} \\
		&=\frac12\lVert pQ_k\rVert^2_{L^2(\Omega\times I)}+s\langle p^2Q_k,v_k\rangle_{L^2(\Omega\times I)}+\frac{s^2}{2}\lVert pv_k\rVert^2_{L^2(\Omega\times I)}+\frac\gamma2\lVert T_k-T_\text{fin}\rVert^2_{L^2(\Omega)}\big|_{t=t_F}+\gamma s\langle T_k-T_\text{fin},Sv_k\rangle_{L^2(\Omega)}\big|_{t=t_F}+s^2\frac\gamma2\lVert Sv_k\rVert^2_{L^2(\Omega)}\big|_{t=t_F} \\
		&=s^2\Big(\frac12\lVert pv_k\rVert^2_{L^2(\Omega\times I)}+\frac\gamma2\lVert Sv_k\rVert^2_{L^2(\Omega)}\big|_{t=t_F}\Big)+s\Big(\langle p^2Q_k,v_k\rangle_{L^2(\Omega\times I)}+\gamma \langle T_k-T_\text{fin},Sv_k\rangle_{L^2(\Omega)}\big|_{t=t_F}\Big)+1\Big(\frac12\lVert pQ_k\rVert^2_{L^2(\Omega\times I)}+\frac\gamma2\lVert T_k-T_\text{fin}\rVert^2_{L^2(\Omega)}\big|_{t=t_F}\Big) \\
		&=s^2\Big(\frac12\lVert pv_k\rVert^2_{L^2(\Omega\times I)}+\frac\gamma2\lVert Sv_k\rVert^2_{L^2(\Omega)}\big|_{t=t_F}\Big)+s\Big(\langle p^2Q_k,v_k\rangle_{L^2(\Omega\times I)}+\gamma \langle S^*(T_k-T_\text{fin}),v_k\rangle_{L^2(\Omega)}\big|_{t=t_F}\Big)+1\Big(\frac12\lVert pQ_k\rVert^2_{L^2(\Omega\times I)}+\frac\gamma2\lVert T_k-T_\text{fin}\rVert^2_{L^2(\Omega)}\big|_{t=t_F}\Big) \\
		&=s^2\Big(\frac12\lVert pv_k\rVert^2_{L^2(\Omega\times I)}+\frac\gamma2\lVert Sv_k\rVert^2_{L^2(\Omega)}\big|_{t=t_F}\Big)+s\Big(\langle p^2Q_k,v_k\rangle_{L^2(\Omega\times I)}+\langle W_k,v_k\rangle_{L^2(\Omega)}\big|_{t=t_F}\Big)+1\Big(\frac12\lVert pQ_k\rVert^2_{L^2(\Omega\times I)}+\frac\gamma2\lVert T_k-T_\text{fin}\rVert^2_{L^2(\Omega)}\big|_{t=t_F}\Big) \\
		&=s^2\Big(\frac12\lVert pv_k\rVert^2_{L^2(\Omega\times I)}+\frac\gamma2\lVert Sv_k\rVert^2_{L^2(\Omega)}\big|_{t=t_F}\Big)-s\lVert v_k\rVert^2_{L^2(\Omega\times I)}+1\Big(\frac12\lVert pQ_k\rVert^2_{L^2(\Omega\times I)}+\frac\gamma2\lVert T_k-T_\text{fin}\rVert^2_{L^2(\Omega)}\big|_{t=t_F}\Big)
	\end{align*}
	This is a quadratic function in $s$ and has the minimum in
	$$s_\text{min}=\frac{\lVert v_k\rVert^2_{L^2(\Omega\times I)}}{\lVert pv_k\rVert^2_{L^2(\Omega\times I)}+\gamma\lVert Sv_k\rVert^2_{L^2(\Omega)}\big|_{t=t_F}}$$
	which requires to solve the (SE) again with input $v_k$ instead of $Q_k$.


	\section{Version 3}
	\begin{align*}
		&\min E(Q,T)=\int\limits_{0}^{t_F}\int\limits_\Omega p(t)Q(x,t)\,\mathrm{d}x\mathrm dt \\
		&\begin{cases}\begin{aligned}
				\rho c\dot T(x,t)&=\nabla\cdot\left(k(x)\nabla T(x,t)\right)+Q(x,t) && \text{in}\ \Omega\times(0,t_F) \\
				-k(x)\partial_\nu T(x,t)&=h(x)\left(T(x,t)-T_\text{out}(x,t)\right) && \text{on}\ \partial\Omega\times[0,t_F]\\
				T(x,0)&=T_\text{ini}(x)&&
		\end{aligned}\end{cases}\tag{SE}\\
		& Q\in U_\text{ad}=\{Q\in L^2(\Omega\times(0,t_F))\,|\,0\leq Q(x,t)\leq Q_\text{max}(x)\ \text{a.e.}\}\\
		&T(x,t_F)\geq T_\text{fin}(x)
	\end{align*}
	The (SE) and thus (WSE) is the same. But we get a new Lagrangian multiplier $\mu\geq0$ and the Lagrangian is
	\begin{align*}
		L(T,q,W,\mu)=\int\limits_{0}^{t_F}\int\limits_\Omega p(t)Q(x,t)\,\mathrm{d}x\mathrm dt+\int\limits_0^{t_F}\int\limits_\Omega \rho c\dot T(x,t)W(x,t)+k(x)\nabla T(x,t)\cdot\nabla W(x,t)-Q(x,t)W(x,t)\,\mathrm dx\mathrm dt+\int\limits_0^{t_F}\int\limits_{\partial\Omega} h(y)\left(T(y,t)-T_\text{out}(y)\right)W(y,t)\,\mathrm dy\mathrm dt+\int\limits_\Omega \mu(x)(T(x,t_F)-T_\text{fin}(x))\mathrm dx
	\end{align*}


	\section{Version 4}
	\begin{align*}
		&\min E(Q_t,T)=\frac12\int\limits_{0}^{t_F}(p(t)Q_t(t))^2\int\limits_\Omega Q_\text{pos}^2(x)\,\mathrm{d}x\mathrm dt +\frac\gamma2\int\limits_\Omega |T(x,t_F)-T_\text{fin}(x)|^2\mathrm dx \\
		&\begin{cases}\begin{aligned}
				\rho c\dot T(x,t)&=\nabla\cdot\left(k(x)\nabla T(x,t)\right)+Q_t(t)Q_\text{pos}(x) && \text{in}\ \Omega\times(0,t_F) \\
				-k(x)\partial_\nu T(x,t)&=h(x)\left(T(x,t)-T_\text{out}(x,t)\right) && \text{on}\ \partial\Omega\times[0,t_F]\\
				T(x,0)&=T_\text{ini}(x)&&
		\end{aligned}\end{cases}\tag{SE}\\
		& Q\in U_\text{ad}=\{Q\in L^2(\Omega\times(0,t_F))\,|\,0\leq Q(x,t)\leq Q_\text{max}(x)\ \text{a.e.}\}\\
		&T(x,t_F)\geq T_\text{fin}(x)
	\end{align*}
	We assume the heat of the heater has the form $Q(x,t)=Q_t(t)Q_\text{pos}(x)$, where $Q_\text{pos}$ is an indicator function describing the (not changing) position of the heater and $Q_t$ describes the power output between $0$ and $Q_\text{max}$. So $Q$ is represented by $Q_t$ describing what setting the thermostat is on, and we can focus just on that function.
	\begin{align*}
		L(Q_t,T,W)&=\frac12\int\limits_{0}^{t_F}(p(t)Q_t(t))^2\int\limits_\Omega Q_\text{pos}^2(x)\,\mathrm{d}x\mathrm dt +\frac\gamma2\int\limits_\Omega |T(x,t_F)-T_\text{fin}(x)|^2\mathrm dx+\int\limits_{0}^{t_F}\int\limits_\Omega \rho c\dot T(x,t)W(x,t)+k(x)\nabla T(x,t)\cdot\nabla W(x,t)\,\mathrm{d}x\mathrm dt+\int\limits_{0}^{t_F}\int\limits_{\partial\Omega} h(y)(T(y,t)-T_\text{out}(y,t))W(y,t)\,\mathrm{d}y\mathrm dt +\int\limits_0^{t_F}Q_t(t)\int\limits_\Omega Q_\text{pos}(x)W(x,t)\,\mathrm dx\mathrm dt \\
		L_T(Q_t,T,W)\psi&=\gamma\int\limits_\Omega (T(x,t_F)-T_\text{fin}(x))\psi(x,t_F)\mathrm dx+\int\limits_{0}^{t_F}\int\limits_\Omega \underbrace{\rho c\dot \psi(x,t)W(x,t)}_{\stapel{&=\left[\int\limits_\Omega \rho c\psi(x,t)W(x,t)\right]_{t=0}^{t=t_F}-\int\limits_0^{t_F}\int\limits_\Omega\rho c\dot W(x,t)\psi(x,t)\,\mathrm dx\mathrm dt \\ &=\int\limits_\Omega \rho c\psi(x,t_F)W(x,t_F)\,\mathrm dx-\int\limits_0^{t_F}\int\limits_\Omega\rho c\dot W(x,t)\psi(x,t)\,\mathrm dx\mathrm dt}}+k(x)\nabla \psi(x,t)\cdot\nabla W(x,t)\,\mathrm{d}x\mathrm dt+\int\limits_{0}^{t_F}\int\limits_{\partial\Omega} h(y)W(y,t)\psi(y,t)\,\mathrm{d}y\mathrm dt \\
		&=\int\limits_\Omega (\gamma(T(x,t_F)-T_\text{fin}(x))+\rho cW(x,t_F))\psi(x,t_F)\mathrm dx+\int\limits_{0}^{t_F}\int\limits_\Omega- \rho c\dot W(x,t)\psi(x,t)+k(x)\nabla \psi(x,t)\cdot\nabla W(x,t)\,\mathrm{d}x\mathrm dt+\int\limits_{0}^{t_F}\int\limits_{\partial\Omega} h(y)W(y,t)\psi(y,t)\,\mathrm{d}y\mathrm dt \tag{WAE} \\
		&=\int\limits_\Omega (\gamma(T(x,t_F)-T_\text{fin}(x))+\rho cW(x,t_F))\psi(x,t_F)\mathrm dx-\int\limits_{0}^{t_F}\int\limits_\Omega \rho c\dot W(x,t)\psi(x,t)+\nabla\cdot(k(x)\nabla W(x,t))\psi(x,t)\,\mathrm{d}x\mathrm dt+\int\limits_{0}^{t_F}\int\limits_{\partial\Omega} h(y)W(y,t)\psi(y,t)+k(x)\partial_\nu W(y,t)\,\mathrm{d}y\mathrm dt \\
		L_{Q_t}(Q_t,T,W)h&=\int\limits_0^{t_F}p^2(t)Q_t(t)h(t)\int\limits_\Omega Q_\text{pos}^2(x)\,\mathrm dx\mathrm dt+\int\limits_0^{t_F}h(t)\int_\Omega Q_\text{pos}(x)W(x,t)\,\mathrm dx\mathrm dt \tag{VI}
	\end{align*}
	Thus, we have the (AE)
	\[\begin{cases}\begin{aligned}
				\rho c\dot W(x,t)&=-\nabla\cdot\left(k(x)\nabla W(x,t)\right) && \text{in}\ \Omega\times(0,t_F) \\
				-k(x)\partial_\nu W(x,t)&=h(x)W(x,t) && \text{on}\ \partial\Omega\times[0,t_F]\\
				W(x,tF)&=-\frac{\gamma}{\rho c}(T(x,tF)-T_\text{fin}(x))&&
		\end{aligned}\end{cases}\tag{SE}
	\]
	and the gradient
	\[\nabla e(Q_t)=p^2(t)Q_t(t)\int\limits_\Omega Q_\text{pos}^2(x)\,\mathrm dx+\int_\Omega Q_\text{pos}(x)W(x,t)\,\mathrm dx\]

	\section{Projection Property}

	\begin{theorem}
		Let $V$ be a linear space (=vector space), $M\subseteq V$ a closed convex set, $f:V\to\mathbb R$ a strictly convex function and consider the unconstraint and constraint minimisation problems
		$$ \min\limits_{x\in V} f(x)\qquad \text{and}\qquad \min\limits_{x\in M} f(x).$$
		Then each problem has at most one solution and if the unconstrained problem has a solution $x_V\in\argmin\limits_{x\in V} f(x)$ then the projection $\pi(x_V)$ is the solution to the constrained problem.
	\end{theorem}
	\begin{proof}
		(i) Since $M$ is arbitrary, proofing that the constrained problem has at most one solution proofs the claim also for the other case by setting $M=V$. Assuming there are two solution $x^1,x^2\in M$, then $(x^1,x^2)\subseteq M$ by convexity of $M$ and for every $x\in(x^1,x^2)$ by strict convexity of $f$ it is $f(x)<\lambda f(x^1)+(1-\lambda)f(x^2)=1\cdot \min\limits_{y\in M} f(y)$, contradicting the minimality of $x^1,x^2$.

		(ii) Assume the minimum over $M$, denoted by $x_M$, would not be the projection. Then $f(x_V)<f(\pi(x_V))$, $f(x_V)<f(x_M)$ by (i) and $f(x_M)<f(\pi(x_V))$ and thus $f(x_V)<f(x_M)<f(\pi(x_V))$. $x_M$ has to be in the boundary of $M$, otherwise we could as in (i) go a little bit into direction of $x_V$ (better: the direction is $x_V-x_M$) and have smaller vallues. Connecting any two points on the trianle $T=\operatorname{span}\{x_V,x_M,\pi(x_V)\}$ we thus get $f(x)<f(\pi(x_V))$ for every $x\in T\setminus\{\pi(x_V)\}$
	\end{proof}


	\section{Remark: minimal time interval}

	If we consider the system without the boundary term, so without heat loss to the outside, then in order to be able to achieve an end temperature close to $T_\text{fin}$ we need enough time, precisely
	$$t=\frac{|\Omega|}{|Q_\text{ad}|}\frac{T_\text{fin}-T_\text{ini}}{\max Q}$$
	much time, which corresponds to heating the room in average to the final temperature using maximum output all the time.

	For small time intervalls we either get a bigger error in the corresponding term of the energy or need a big $\max Q$. Looking at the Theta-method with $Y_n$ being any solution (of SE or AE) for only a given time $t_n$ the next solution for $t_{n+1}$ with the help of the right hand side extracted by looking at the time change of the system will be
	$$Y_{n+1}=Y_n+\Delta t\left(\theta F(t_n,Y_n)+(1-\theta)F(t_{n+1},Y_{n+1})\right)$$
	which gets solved by the Newton method iterating over $k$
	$$Y_{n+1}^{k+1}=Y_{n+1}^k-\left(G'(Y_{n+1}^k)\right)^{-1}G(Y_{n+1}^k)$$
	where
	$$G(Y)=Y-\Delta t\left(\theta F(t_n,Y_n)+(1-\theta)F(t+h,Y)\right)-Y_n,\qquad\qquad G'(Y)=I-(1-\theta)\Delta t\partial_Y F(t_{n+1},Y)$$


	\section{Adjoint}

	\begin{theorem}
		The Adjoint solution operator from version 2 satisfies
		$$\langle S^*(T-T_\text{fin}),h\rangle = \langle T-T_\text{fin},S(h)\rangle$$
	\end{theorem}
	\begin{proof}
		We remind about the weak forms (WSE) with test function $W$ and (WAE) (from $L_T$ before using Green) with test function $T$
		\begin{align}
			\int\limits_0^{t_F}\int\limits_\Omega \rho c\dot T(x,t)W(x,t)\,\mathrm dx\mathrm dt=\int\limits_0^{t_F}\int\limits_\Omega -k(x)\nabla T(x,t)\cdot\nabla W(x,t)+Q(x,t)W(x,t)\,\mathrm dx\mathrm dt-\int\limits_0^{t_F}\int\limits_{\partial\Omega} h(y)\left(T(y,t)-T_\text{out}(y)\right)W(y,t)\,\mathrm dy\mathrm dt \tag{WSE} \\
			\int\limits_0^{t_F}\int\limits_\Omega \rho c\dot T(x,t)W(x,t)\,\mathrm dx\mathrm dt=-\int\limits_0^{t_F}\int\limits_\Omega k(x)\nabla T(x,t)\cdot\nabla W(x,t)\,\mathrm dx\mathrm dt-\gamma\int\limits_\Omega \left(T(x,t_F)-T_\text{fin}(x)\right)T(x,t_F)\mathrm dx-\int\limits_0^{t_F}\int\limits_{\partial\Omega} h(y)W(y,t)T(y,t)\,\mathrm dy\mathrm dt \tag{WAE}
		\end{align}
		Subtracting (WAE) from (WSE) we get
		\begin{align*}
			0=\int\limits_0^{t_F}\int\limits_\Omega Q(x,t)W(x,t)\,\mathrm dx\mathrm dt+\int\limits_0^{t_F}\int\limits_{\partial\Omega} h(y)T_\text{out}(y)W(y,t)\,\mathrm dy\mathrm dt+\gamma\int\limits_\Omega \left(T(x,t_F)-T_\text{fin}(x)\right)T(x,t_F)\mathrm dx
		\end{align*}
		With these it is
		\begin{align*}
			\langle W,S(Q)\rangle&=\int\limits_0^{t_F}\int\limits_\Omega W(x,t)T(x,t)\,\mathrm dx\mathrm dt \\
			&=\int\limits_0^{t_F}\int\limits_\Omega W(x,t)\left(T_\text{ini}(x)+\int\limits_0^{t} \dot T(x,s)\,\mathrm ds\right)\,\mathrm dx\mathrm dt \\
			&=\int\limits_0^{t_F}\int\limits_\Omega W(x,t)T_\text{ini}(x)+\int\limits_0^{t} \dot T(x,s)W(x,t)\,\mathrm ds\,\mathrm dx\mathrm dt \\
			&=\int\limits_0^{t_F}\int\limits_\Omega W(x,t)T_\text{ini}(x)+\int\limits_0^{t} \dot T(x,s)W(x,t)\,\mathrm ds\,\mathrm dx\mathrm dt \\
			&=
		\end{align*}
	\end{proof}
	In general it looks like this
	\begin{theorem}
		Assume the problem is
		\[\min E(Q,T)\quad\text{such that}\quad N(Q,T)=0\]
		where $N$ contains only linear Terms of $Q$ and $T$.
	\end{theorem}
	\begin{proof}
		The Lagrangian is
		\[L(Q,T,W)=E(Q,T)+WN(Q,T)\]
		and the adjoint equation is
		\[L_T(Q,T,W)h=E_T(Q,T)+WN_T(Q,T)h\]

	\end{proof}


	\section{Pontryagin's minimum (maximum) principle}
	See \cite{Jaksch2017ImplementationadjointthermalsolverinverseproblemsconferencePaper}.
	We can also get the optimal adjoint by Pontryagin's minimum principle, which may be used here to even out some holes in our argumentation. We assume version 2.

	\begin{theorem}[Pontryagin's minimum principle (version \cite{Jaksch2017ImplementationadjointthermalsolverinverseproblemsconferencePaper})]
		Consider the problem
		\begin{align*}
			\min\limits_{q\in Q} \Phi(y(t_f))+\int\limits_0^{t_f} F\big(y(t),q(t)\big)\,\mathrm dt\\
			\dot y=f\big(y(t),q(t)\big),\quad y(0)=y_0
		\end{align*}
		Let the Hamiltonian be given by $H(q,y,p,t)=F\big(y(t),q(t)\big)+p(t)f\big(y(t),q(t)\big)$. Then the optimal adjoint solves
		\[\dot p=H_y\qquad p(t_f)=\Phi_y(y(t_f))\]
	\end{theorem}
	Identifying the functions for our problem and using $y=T$, $q=Q$, $p=W$ we get (notice that here we write everything in the strong version)
	\begin{align*}
		\Phi(T(t_f))&= \int\limits_\Omega \big(T(x,t_f)-T_\text{fin}(x)\big)^2\,\mathrm dx\\
		F\big(T(t),Q(t)\big)&= \int\limits_\Omega (pQ)^2\,\mathrm dx\\
		f\big(T(t),Q(t)\big)&= \frac{1}{\rho c}\left(\nabla\cdot\left(k(x)\nabla T(x,t)\right)+Q(x,t)\right)\\%-\int\limits_{\partial\Omega} \frac{h(y)}{\rho c}\left(T(y,t)-T_\text{out}(y)\right)\,\mathrm dy \\
		H(Q,T,W,t)&= \int\limits_\Omega (pQ)^2\,\mathrm dx-\int\limits_0^{t_F}\int\limits_\Omega k(x)\nabla T(x,t)\cdot\nabla W(x,t)-Q(x,t)W(x,t)\,\mathrm dx\mathrm dt-\int\limits_0^{t_F}\int\limits_{\partial\Omega} h(y)\left(T(y,t)-T_\text{out}(y)\right)W(y,t)\,\mathrm dy\mathrm dt
	\end{align*}
	and thus the adjoint satisfies
	\[\dot W(x,t)h=\]


	\section{Other}
	\begin{itemize}
		\item Remember the Courant-Friedrichs-Lewy (CFL) condition: If the numerical solution is screwed up it might be because space step length / time step length is too small. Honestly, I do not know the specific condition for that problem, it might even be reversed here, but this should be considered.
		\item Galerkin method or Ritz method - mention them and I think we use Galerkin
		\item Pontryagin's maximum principle??? Gives the adjoint of a system like ours, but it might be hard to transform
	\end{itemize}

	\section{Remarks Henning}

	\cite{TabarkiMabrouk2012ModelingEperimentationHeatExchangesRoomHeatedRadiatorjournalArticle} Introduction/motivation when saying the problem has been studied by others, conclusion when saying what other models can be studied (they study a 3D model with radiation term and consider a Navier-Stokes model)

	Either for the adjoint or the Variation you could state that this is the sensitivity of E with respect to T.

	Heat conduction in room/problem statement and justification: \cite{HahnOzisik2012Heatconductionbook}

	Gridap citing or when talking about software: \cite{BadiaVerdugo2020GridapextensibleFiniteElementtoolboxJuliajournalArticle} \cite{VerdugoBadia2022softwaredesignGridapFiniteElementpackagebasedJuliaJITcompilerjournalArticle}

	Pontryafin's minimum principle, a direct solution method \cite{VanKeulen2014IndirectsolutionoptimalcontrolproblemspurestateconstraintjournalArticle}

	\cite{TrapnesOptimalTemperatureControlRoomsjournalArticle} Introduction/motivation when saying the problem has been studied by others, model, conclusion when talking about other models

	\cite{PulsipherZhangHongistoZavala2022UnifyingModelingAbstractionInfiniteDimensionalOptimizationjournalArticle} other numerical methods

	\cite{LoosdirekteMethodefuerOptimalsteuerungsproblememitparabolischenDifferentialgleichungenwebpage} weak formulation, numerical method (discretisation, semi-discretisation, FE method), when saying that we use a direct method, other methods (IPOPT), Lagrangian

	\cite{LewisNithiarasuSeetharamu2004FundamentalsFiniteElementMethodHeatFluidFlowbook} model (chapter 6), parameters, in conclusion when talking about other models, FE, Lagrangian, Rayleigh-Ritz + Galerkin

	\cite{Jaksch2017ImplementationadjointthermalsolverinverseproblemsconferencePaper} model, Lagrangian, Variation of Lagrangian, numerical method with gradient

	\cite{IncroperaDewittBergmanLavine2007Fundamentalsheatmasstransferbook} model (esp chapter 5), parameters, conclusion other models (include radiation, Stefan-Boltzmann constant)

	\cite{HuangSeidelJiaPaschkeBraunig2021EnergyOptimalControlMultivalentBuildingEnergySystemusingMachineLearningconferencePaper} other numerical method, slightly other model including solar transfer

	\cite{Hinze2005VariationalDiscretizationConceptControlConstrainedOptimizationLinearQuadraticCasejournalArticle} unique solution of quadratic problem, numerical solution has error in O($h^2$) using discretisation and weak form, our gradient method, active set method

	\cite{HahnOzisik2012Heatconductionbook} model, analytical solution

	\cite{GustafssonRonnqvist2008OptimalheatinglargeblockflatsjournalArticle} extended problem for multiple flats, simplex method like

	\cite{GongHinzeZhou2014finiteelementmethodDirichletboundarycontrolproblemsgovernedparabolicPDEspreprint} model, error bound in equation (1.4), weak form, adjoint, regularity of solution, FE

	\cite{GockenbachSchmidtke2009NewtonslawheatingheatequationjournalArticle} model, FE

	\cite{FilipovHristovAvdzhievaFarago2023CoupledPDEODEModelNonlinearTransientHeatTransferConvectionHeatingBoundaryNumericalSolutionImplicitTimeDiscretizationSequentialDecouplingjournalArticle} model + advanced nonlinear model (for conclusion), finite difference method

	\cite{Evans2010Partialdifferentialequationsbook} for everything about PDEs (sobolev spaces, weak formulationheat equation, solution theorem)

	\cite{BiralBertolazziBosetti2016NotesNumericalMethodsSolvingOptimalControlProblemsjournalArticle} model, Pontryagin's minimum principle,, Lagrangian (or here: Hamiltonian), indirect method

	\cite{BauschkeCombettes2017ConvexAnalysisMonotoneOperatorTheoryHilbertSpacesbook} general convex optimisation theory, spaces, projection

	\cite{Alt2016LinearFunctionalAnalysisbook} sobolev spaces, differentiability ... general stuff

	\cite{Mordukhovich2006VariationalAnalysisGeneralizedDifferentiationbook} Variational analysis (weak formulation is good), differentiation in banach spaces, optimisation with variation

	\cite{Luenberger1969OptimizationVectorSpaceMethodsbook} adjoint operator, generalised differentiability, Pontryagin's maximum principle (chap 9.6), semi-discretisised problem (here: solve dynamical systems after that)

	\cite{Troeltzsch2009OptimaleSteuerungpartiellerDifferentialgleichungenTheorieVerfahrenundAnwendungenbook} Tröltzsch ... you know whats in there (this is the german version)

	\cite{Liberzon2012CalculusVariationsOptimalControlTheoryConciseIntroductionbook} variational analysis (weak formulation is good), fundamental theorem of variational analysis, general optimal control, minimum principle

	\cite{LoosdirekteMethodefuerOptimalsteuerungsproblememitparabolischenDifferentialgleichungenwebpage} model, weak formulation, FE, Lagrangian

	\cite{Stengel1994OptimalControlEstimationbook} general problems, Lagrangian, good choices for gamma (page 196 ff.), optimality conditions

	\printbibliography
\end{document}